\section{Related Work}

%\wl{1) tweaking pixels/words (I guess this is not the novelty). 2) Multi-view (compare with Contrastive multiview coding). 3) Nearest neighbor}
% \checkme{more related work of NN}
\subsection{Pre-trained Models}
The critical idea of pre-training is to first extract general knowledge implicitly from the massive amount of data and then transfer the knowledge to versatile downstream tasks~\citep{han2021pre}. 
Big NLP models~\citep{devlin2018bert, brown2020language} yield unprecedented performance via learning from tremendous language data over the Internet and labor-free supervision within the language itself.
% With the transformer architecture~\citep{vaswani2017attention} and tremendous language data over the Internet, big pre-trained NLP models, such as BERT-series~\citep{devlin2018bert, liu2019roberta} and GPT-series~\citep{radford2019language, brown2020language}, have dominated this area~\citep{han2021pre}.
% Besides the exquisite architecture, the great success comes from the tremendous language data over the Internet and labor-free supervision~\citep{yang2019xlnet, radford2019language, devlin2018bert} within the language itself.
% Benefiting from task-agnostic objectives such as auto-regressive~\citep{yang2019xlnet, radford2019language} and masked language modeling~\citep{devlin2018bert}, we can directly leverage the unlabeled Internet text data and improve the accuracy via scaling up data, model, and compute~\citep{brown2020language, fedus2021switch}.
In the field of CV, supervised pre-training on ImageNet is still the standard practice. 
While achieving great success on downstream CV tasks~\citep{girshick2014rich, long2015fully, vinyals2015show} , this supervised manner is hard to scale.
% supervised pre-training on a well labeled dataset, like ImageNet~\citep{deng2009imagenet}, then fine-tuning on labeled target dataset is still the standard practice. 
% This paradigm needs a pre-defined fixed set of object categories (e.g. 1000 in ImageNet) and tremendous human labeling effort, thus greatly limiting its generality and scalablity. 
To address this challenge, our \workname learns directly from image-text pairs that are abundant across the Internet. More importantly, by exploiting the widespread supervision within the pairs, our \workname is more data-efficient than the prior art.
